
\documentclass[11pt,a4paper]{article}

%% ============================================================
%% MTPI Core Specification v1.0 - LaTeX Formalization
%% Meta-Theorem of Prime Identity: A Proof-First Framework
%% for Verifiable Computation
%% Author: Ryan O. Van Gelder
%% Affiliation: Citizen Gardens Institute for Mathematical Discovery
%% License: CC-BY-4.0 International
%% Date: January 2026
%% DOI: 10.5281/zenodo.18248116
%% ============================================================

\usepackage[utf8]{inputenc}
\usepackage[T1]{fontenc}
\usepackage{amsmath,amssymb,amsthm}
\usepackage{mathtools}
\usepackage{hyperref}
\usepackage{listings}
\usepackage{xcolor}
\usepackage{booktabs}
\usepackage{enumitem}
\usepackage{geometry}
\usepackage{fancyhdr}
\usepackage{titlesec}

\geometry{margin=1in}

%% Theorem environments
\newtheorem{theorem}{Theorem}[section]
\newtheorem{lemma}[theorem]{Lemma}
\newtheorem{corollary}[theorem]{Corollary}
\newtheorem{definition}[theorem]{Definition}
\newtheorem{proposition}[theorem]{Proposition}
\newtheorem{remark}[theorem]{Remark}

%% Custom commands
\newcommand{\MTPI}{\textsf{MTPI}}
\newcommand{\PIRTM}{\textsf{PIRTM}}
\newcommand{\CSL}{\textsf{CSL}}
\newcommand{\Archivum}{\textsf{Archivum}}
\newcommand{\Tensor}[1]{\mathcal{T}_{#1}}
\newcommand{\EthicalField}{\mathcal{E}_t}
\newcommand{\SovereigntyTensor}{\sigma_i(t)}
\newcommand{\Commutator}[2]{[#1, #2]}
\newcommand{\PrimeSet}[1]{\mathbb{P}(N_{#1})}
\newcommand{\Poseidon}{\textsc{Poseidon}}
\newcommand{\MillerRabin}{\textsc{Miller-Rabin}}
\newcommand{\XiOp}{\Xi(t)}

%% Code listing style
\definecolor{codegreen}{rgb}{0,0.6,0}
\definecolor{codegray}{rgb}{0.5,0.5,0.5}
\definecolor{codepurple}{rgb}{0.58,0,0.82}
\definecolor{backcolour}{rgb}{0.95,0.95,0.92}

\lstdefinestyle{mystyle}{
    backgroundcolor=\color{backcolour},   
    commentstyle=\color{codegreen},
    keywordstyle=\color{magenta},
    numberstyle=\tiny\color{codegray},
    stringstyle=\color{codepurple},
    basicstyle=\ttfamily\footnotesize,
    breakatwhitespace=false,         
    breaklines=true,                 
    captionpos=b,                    
    keepspaces=true,                 
    numbers=left,                    
    numbersep=5pt,                  
    showspaces=false,                
    showstringspaces=false,
    showtabs=false,                  
    tabsize=2
}
\lstset{style=mystyle}

%% Header/Footer
\pagestyle{fancy}
\fancyhf{}
\rhead{MTPI Core Specification v1.0}
\lhead{Defensive Publication}
\rfoot{Page \thepage}
\lfoot{CC-BY-4.0}

%% ============================================================
\begin{document}

\begin{titlepage}
    \centering
    \vspace*{2cm}

    {\Huge\bfseries Meta-Theorem of Prime Identity\\[0.5cm]}
    {\LARGE\itshape A Proof-First Framework for Verifiable Computation\\[1cm]}

    {\Large\textbf{MTPI Core Specification v1.0}\\[2cm]}

    {\large
    Ryan O. Van Gelder\\[0.3cm]
    \textit{Citizen Gardens Institute for Mathematical Discovery}\\
    Worthington, Ohio, USA\\[2cm]
    }

    {\large Publication Date: January 2026\\[0.5cm]}

    \vfill

    {\normalsize
    \textbf{Document Classification:} Defensive Publication / Technical Specification\\
    \textbf{License:} Creative Commons Attribution 4.0 International (CC-BY-4.0)\\
    \textbf{Version:} 1.0.0\\[0.3cm]
    \textbf{DOI:} \href{https://doi.org/10.5281/zenodo.18248116}{10.5281/zenodo.18248116}
    }
\end{titlepage}

%% ============================================================
\tableofcontents
\newpage

%% ============================================================
\begin{abstract}
The \textbf{Meta-Theorem of Prime Identity (MTPI)} is an architectural pattern and protocol stack for building \emph{proof-first} digital systems. Instead of relying on institutional trust, surveillance, or opaque algorithms, MTPI requires that every state transition be accompanied by a cryptographic proof that the transition is \emph{lawful}, \emph{ethically constrained}, and \emph{user-sovereign}.

Lawfulness is defined not only in terms of syntactic validity (e.g., input formats, signatures) but also with respect to higher-order invariants: ethical constraints, bounded drift, jurisdictional rules, and user consent.

\textbf{Core Claim:} Every lawful identity in the system can be represented as a composition of prime identities, and this decomposition is stable under lawful recursion.

An MTPI identity is anchored by prime-based checks and is mathematically constrained by strict time-drift bounds, typically $\delta(t) \leq 0.3$. This ensures that proofs are not only valid but also \emph{live}, preventing replay attacks and impersonation by ensuring a proof is linked to a specific, verifiable moment in time.
\end{abstract}

%% ============================================================
\section{Mathematical Foundation}

\subsection{Prime-Indexed Recursive Tensor Mathematics (PIRTM)}

\begin{definition}[Recursive Tensor Evolution]
The recursive tensor dynamics within PIRTM are defined by the evolution equation:
\begin{equation}
    \Tensor{t+1} = \lambda \sum_{p_i \in \PrimeSet{t}} \Tensor{t} + \mathcal{F}_t
    \label{eq:tensor_evolution}
\end{equation}
where:
\begin{itemize}[noitemsep]
    \item $\lambda$ is the \textbf{Universal Multiplicity Constant} (cognitive stabilizer), with $\lambda \geq 1$
    \item $\PrimeSet{t}$ is the set of primes up to index $N_t$
    \item $\mathcal{F}_t$ represents forcing inputs (cognitive, semantic, or external)
    \item The condition $\lambda \geq 1$ ensures exponential convergence
\end{itemize}
\end{definition}

\subsection{Contractive Dynamics Guarantee}

\begin{theorem}[Stability Guarantee]
The system's reliability is underpinned by a foundational mathematical guarantee of stability derived from the principle of contractive dynamics, governed by the core inequality:
\begin{equation}
    \|\Phi_t\| \leq q^t, \quad q < 1
    \label{eq:contractive}
\end{equation}
This ensures that with every computational step, the distance between any two possible computational paths is guaranteed to shrink.
\end{theorem}

\begin{proof}
The proof of convergence for the recursive operator $\Phi_t$ follows from the Banach Fixed-Point Theorem. Given the contractive mapping property with contraction constant $q < 1$, the sequence of iterates converges to a unique fixed point in the complete metric space of tensor states.
\end{proof}

\subsection{Recursive Operator Evolution}

\begin{definition}[Operator Dynamics]
The recursive operator evolves under the differential relation:
\begin{equation}
    \frac{d\Phi_t}{dt} = \mathcal{M}\Phi_t - \Commutator{\mathcal{M}}{\Phi_t}
    \label{eq:operator_evolution}
\end{equation}
where $\mathcal{M}$ is a Hilbert-Schmidt operator and $\Commutator{\mathcal{M}}{\Phi_t} := \mathcal{M}\Phi_t - \Phi_t\mathcal{M}$ denotes the commutator. This formulation follows the dynamic recursive operator framework established in \cite{vanosdol2025}.
\end{definition}

\begin{remark}
The commutator $\Commutator{A}{B}$ measures the degree to which operators $A$ and $B$ fail to commute. When $\Commutator{A}{B} = 0$, the operators are said to commute, implying that the order of application does not affect the outcome.
\end{remark}

\begin{definition}[Dynamic Recursive Operator]
The dynamic recursive operator $\XiOp$ governs tensor evolution in PIRTM systems through $\Lambda_m$-indexed contraction mappings. Stability in DRMM/PIRTM is achieved via entangled feedback loops that preserve prime-indexed structure under recursive application \cite{vanosdol2025}.
\end{definition}

\subsection{Related Work and Prior Art}

MTPI extends foundational cryptographic and number-theoretic primitives:

\begin{itemize}[noitemsep]
    \item \textbf{Groth16} (2016): Foundational ZK-SNARK construction \cite{groth2016}
    \item \textbf{Poseidon Hash} (2021): Algebraic hash function optimized for ZK circuits \cite{grassi2021}
    \item \textbf{Miller-Rabin} (1980): Probabilistic primality testing \cite{rabin1980}
    \item \textbf{Van Osdol} (2025): Dynamic recursive operator $\XiOp$ and PIRTM stability analysis via Hilbert-Schmidt dynamics and prime-indexed contraction mappings \cite{vanosdol2025}
\end{itemize}

MTPI introduces \emph{prime-indexed identity gating}, \emph{drift-bounded state evolution}, and \emph{silence-by-default semantics} as novel extensions.

%% ============================================================
\section{Prime-Indexed Identity}

\subsection{Genesis State and Prime Index}

\begin{definition}[Identity Anchoring]
Every governed entity (user, device, process, model, contract) is anchored in a \textbf{genesis state} and a \textbf{prime index}:
\begin{itemize}[noitemsep]
    \item The genesis state encodes initial conditions (e.g., enrollment seed, device attestation, initial model snapshot)
    \item The prime index ties that genesis to a unique slot within the system's identity lattice
\end{itemize}
\end{definition}

\begin{definition}[Identity Computation]
The identity is computed via a \Poseidon{} hash of the genesis state and local salt, then mapped to a prime index within a large, sparse space:
\begin{equation}
    \text{Identity} = \Poseidon(0\text{x}39\text{e}0, \text{salt})
    \label{eq:identity}
\end{equation}
Validated via \MillerRabin{} primality testing.
\end{definition}

\subsection{Prime-Gated Activation}

State transitions are gated by prime-based checks, where actions are only permissible during specific \emph{prime epochs}.

\begin{lstlisting}[language=C,caption={Prime-Gated Activation (Solidity pseudocode)}]
function activatePrimeGate(uint256 primeIndex, bytes32 proofHash) external {
    if (!isCertifiedPrime(primeIndex)) 
        revert InvalidPrime(primeIndex);
    if (primeIndex <= state.lastPrimeIndex) 
        revert PrimeOrderViolation(primeIndex, state.lastPrimeIndex + 1);
    // ... activation logic
}
\end{lstlisting}

\subsection{Drift Bounds}

\begin{definition}[Drift Constraint]
Drift measures the distance between an identity's current state and its genesis or prior lawful checkpoints. The system enforces:
\begin{equation}
    \delta(t) \leq 0.3
    \label{eq:drift_bound}
\end{equation}
\end{definition}

\begin{proposition}[Drift Enforcement]
When drift exceeds the threshold, the system must either:
\begin{enumerate}[noitemsep]
    \item Refuse further transitions until recertification occurs, or
    \item Restrict capabilities and enter constrained mode
\end{enumerate}
\end{proposition}

%% ============================================================
\section{Conscious Sovereignty Layer (CSL)}

\subsection{Ethical Tensor Field}

\begin{definition}[Conscious Sovereignty Layer]
The \CSL{} is a governance and policy layer that encodes:
\begin{itemize}[noitemsep]
    \item \textbf{Ethical constraints:} beneficence (do good), non-maleficence (do no harm), respect for autonomy
    \item \textbf{Sovereignty rules:} which jurisdictions apply to a given identity or action
    \item \textbf{Silence Clause semantics:} specifying when systems must default to inaction
\end{itemize}
\end{definition}

\subsection{CSL Commutation Relation}

\begin{theorem}[Admissibility Criterion]
A transition is admissible only if, when interpreted under CSL, it commutes with the ethical operators:
\begin{equation}
    \Commutator{\mathcal{M}}{\EthicalField} = 0 \quad \Longrightarrow \quad \mathcal{M} \text{ admissible}
    \label{eq:csl_commutation}
\end{equation}
where $\EthicalField$ is the Ethical Tensor Field encoding conserved invariants.
\end{theorem}

\subsection{Sovereignty Tensor}

\begin{definition}[Sovereignty Encoding]
The Sovereignty Tensor $\SovereigntyTensor$ encodes agent permissions at node $i$ and time $t$:
\begin{equation}
    \SovereigntyTensor \in \{0, 1\}^n
    \label{eq:sovereignty_tensor}
\end{equation}
where each binary element corresponds to constraints: consent, security, context alignment.
\end{definition}

\subsection{Recursive Opt-Out (Silence Clause)}

\begin{theorem}[Silence Clause]
The rule guaranteeing an agent's right to disengage:
\begin{equation}
    \Tensor{t+1}^{(i)} = \Tensor{t}^{(i)} \quad \text{if} \quad \SovereigntyTensor = \mathbf{0}
    \label{eq:silence_clause}
\end{equation}
This mathematically freezes the agent's state upon opt-out.
\end{theorem}

\begin{corollary}[Default Behavior]
In ambiguous situations or in the absence of a valid proof, the system's default behavior is \emph{silence}---to take no action.
\end{corollary}

%% ============================================================
\section{Archivum Audit Schema}

\subsection{$\Lambda$-Trace Structure}

\begin{definition}[$\Lambda$-Trace]
$\Lambda$-Trace serves as the canonical decision and audit layer, binding off-chain decisions to on-chain assets via a single unique hash---the \texttt{ltraceHash}. A $\Lambda$-Trace object captures the full context of a decision:
\begin{itemize}[noitemsep]
    \item Scores and evaluation metrics
    \item Ethics verdict (CSL compliance)
    \item Link to supporting evidence (IPFS/decentralized storage)
\end{itemize}
\end{definition}

\begin{lstlisting}[language=C,caption={$\Lambda$-Trace Interface (TypeScript)}]
interface LambdaTrace {
    primeId: string;       // Prime-indexed identity hash
    proofHash: bytes32;    // ZK proof hash
    vkHash: bytes32;       // Verification key hash
    R96: number;           // Resonance budget
    delta: number;         // Drift metric
    timestamp: uint256;
    lawful: boolean;
}
\end{lstlisting}

\subsection{Provenance Ledger}

\begin{definition}[\Archivum{}]
\Archivum{} is an append-only, prime-indexed store that records:
\begin{itemize}[noitemsep]
    \item Hashes of zero-knowledge proofs
    \item Minimal public metadata (timestamps, drift metrics)
    \item Prime identity references
    \item Resonance/participation weights
\end{itemize}
\end{definition}

\begin{remark}[Privacy Guarantee]
\textbf{Critical:} \Archivum{} does not store raw personal data. It enables regulator-grade auditability without surveillance.
\end{remark}

\subsection{Post-Quantum Integrity}

\begin{theorem}[Hash Chain Immutability]
Immutability is guaranteed by a recursive state hash chain:
\begin{equation}
    H_{n+1} = \mathcal{H}(B_n \| H_n \| Q_{\text{ent}}(n))
    \label{eq:hash_chain}
\end{equation}
where $Q_{\text{ent}}(n)$ is a high-entropy source (e.g., quantum random number generator or cryptographic entropy pool), providing resistance against future quantum threats.
\end{theorem}

%% ============================================================
\section{Conformance Requirements}

\subsection{Technical Invariants}

All MTPI-conformant implementations must satisfy:

\begin{table}[h]
\centering
\begin{tabular}{@{}lll@{}}
\toprule
\textbf{Invariant} & \textbf{Requirement} & \textbf{Enforcement} \\ \midrule
Prime-Gated Activation & $\text{isPrime}(\text{epochIndex}) = \text{true}$ & Miller-Rabin circuit \\
Drift Bound & $\delta(t) \leq 0.3$ & Client-side + on-chain verification \\
CSL Commutation & $[\mathcal{M}, \mathcal{E}_t] = 0$ & Proof must demonstrate commutation \\
Replay Protection & $\text{usedProofHashes}[\text{proofHash}] = \text{false}$ & Nullifier mapping \\
Silence-by-Default & No action without valid proof & Revert on verification failure \\ \bottomrule
\end{tabular}
\caption{Technical invariants for MTPI conformance}
\label{tab:invariants}
\end{table}

\subsection{Surveillance Fork Definition}

\begin{definition}[Surveillance Fork]
An implementation constitutes a \textbf{Surveillance Fork} if it:
\begin{enumerate}[noitemsep]
    \item \textbf{EMITS} any of the following WITHOUT an on-device consent artifact:
    \begin{itemize}[noitemsep]
        \item Device fingerprints (canvas, WebGL, audio context)
        \item Persistent identifiers linkable across $>1$ session
        \item Location data at precision $\leq$ city level
        \item Biometric templates or derivatives
    \end{itemize}
    \item \textbf{OR FAILS} the Emission Gating Test
    \item \textbf{OR STORES} on remote infrastructure raw user content without client-side encryption where user does not hold decryption key
\end{enumerate}
\end{definition}

\subsection{Emission Gating Test}

\begin{lstlisting}[language=bash,caption={Emission Gating Test}]
#!/bin/bash
# Emission Gating Test for MTPI Conformance
# Requires containerization/sandboxing (Docker or equivalent)

docker run --network=none $IMPLEMENTATION_IMAGE sleep 60 &
tcpdump -i any -c 1 -w /tmp/capture.pcap &
# Enable network
docker network connect bridge $CONTAINER_ID
# Wait for idle
sleep 300
# Analyze capture
if pcap_has_unlisted_endpoints /tmp/capture.pcap $MANIFEST; then
    echo "FAIL: Surveillance Fork detected"
    exit 1
fi
echo "PASS: Emission Gating Test"
\end{lstlisting}

\begin{remark}
For non-containerized implementations, use \texttt{strace}/\texttt{dtrace} syscall tracing or equivalent sandboxing.
\end{remark}

\subsection{Proof Lifecycle}

The end-to-end lifecycle of an MTPI transaction:

\begin{enumerate}[noitemsep]
    \item \textbf{Local Proof Generation:} Client uses circuit (WASM binary) and proving key to generate Groth16 ZK-proof
    \item \textbf{Local Verification:} Client verifies proof against local verification key; if fail, halt
    \item \textbf{Policy Oracle Evaluation:} Client evaluates Prime Gate + Drift Gate; verdict: submit or silent
    \item \textbf{On-Chain Submission:} If submit, proof and minimal public signals sent to verifier contract
    \item \textbf{State Transition:} Upon verification, contract executes transition, records proof hash as nullifier
    \item \textbf{Audit Event Emission:} Contract emits privacy-preserving audit event (hashes only)
\end{enumerate}

%% ============================================================
\section{Reference Implementation}

\subsection{Core Contract Interface}

\begin{lstlisting}[language=C,caption={IMTPICore Interface (Solidity)}]
// SPDX-License-Identifier: LProof-License-v1.2
pragma solidity ^0.8.20;

interface IMTPICore {
    function quantumTransition(
        uint256[2] calldata a,
        uint256[2][2] calldata b,
        uint256[2] calldata c,
        uint256[1] calldata input,
        bytes32 proofHash
    ) external returns (bytes32 newStateHash);

    function enterSilentRecovery(bytes32 proofHash) external;
    function getCurrentState() external view returns (bytes32);
    function checkPrimeAccess(address entity, uint256 primeIndex) 
        external view returns (bool);
}
\end{lstlisting}

\subsection{Circom Circuit Template}

\begin{lstlisting}[language=C,caption={RootContract Circuit (Circom 2.1+)}]
pragma circom 2.0.0;

include "poseidon.circom";
include "millerrabin.circom";
include "comparators.circom";

template RootContract() {
    signal input identityHash;
    signal input stateCommit;
    signal input primeIndex;
    signal input drift;

    signal output newStateHash;
    signal output lawful;

    // Prime-gate check
    component isPrime = MillerRabin();
    isPrime.n <== primeIndex;
    isPrime.result === 1;

    // Drift bound check (Circom 2.1 compliant)
    // 0.3 represented as 3e17 in fixed-point (scale 1e18)
    signal driftBound;
    driftBound <== 300000000000000000;

    component driftCheck = LessThan(64);
    driftCheck.in[0] <== drift;
    driftCheck.in[1] <== driftBound;
    signal driftValid;
    driftValid <== driftCheck.out;
    driftValid === 1;

    // State evolution via Poseidon
    component hasher = Poseidon(3);
    hasher.inputs[0] <== identityHash;
    hasher.inputs[1] <== stateCommit;
    hasher.inputs[2] <== primeIndex;

    newStateHash <== hasher.out;
    lawful <== 1;
}

component main = RootContract();
\end{lstlisting}

%% ============================================================
\section{Intellectual Property Declaration}

\subsection{Public Domain Designation (as of v1.0.0)}

The following concepts are irrevocably placed in the public domain and may not be patented, enclosed, or restricted, effective as of specification version \textbf{v1.0.0}:

\begin{enumerate}[noitemsep]
    \item \textbf{MTPI} (Meta-Theorem of Prime Identity) mathematical formulation
    \item \textbf{PIRTM} (Prime-Indexed Recursive Tensor Mathematics)
    \item \textbf{PETC} (Prime-Eigenmode Tensor Calculus)
    \item \textbf{DRMM} (Recursive Decodability and Multiplicity Mathematics)
    \item \textbf{CSL} (Conscious Sovereignty Layer) ethics framework, including Silence Clause
    \item Prime-gated invariant specifications
    \item RootContract interface specifications
    \item \Archivum{} audit logic (specification-level)
    \item Dynamic recursive operator $\XiOp$ framework and $\Lambda_m$-indexed contraction mappings
\end{enumerate}

\begin{remark}[Version Scope]
This public domain designation applies to specification versions \texttt{1.x.x}. Future major versions (\texttt{2.0.0}+) may extend but not retract public domain scope.
\end{remark}

\subsection{Guiding Principle}

\begin{quote}
\textit{The Constitution is open; the tools may be owned.}
\end{quote}

This bifurcated strategy ensures no single entity can control the foundational layer while enabling commercialization of high-value implementations built atop the open foundation.

%% ============================================================
\section{Alternative Embodiments}

\subsection{Non-Blockchain Deployments}

MTPI may be implemented without blockchain anchoring, using timestamped append-only logs (e.g., Certificate Transparency, Trillian) for \Archivum{}.

\subsection{Alternative Proof Systems}

Groth16 may be substituted with:
\begin{itemize}[noitemsep]
    \item \textbf{PLONK:} Universal trusted setup
    \item \textbf{Halo2:} No trusted setup
    \item \textbf{STARKs:} Post-quantum, larger proofs
\end{itemize}

\subsection{Alternative Hash Functions}

\Poseidon{} may be substituted with:
\begin{itemize}[noitemsep]
    \item \textbf{Rescue:} Similar algebraic structure
    \item \textbf{Blake3:} Non-algebraic, requires different circuit design
\end{itemize}

\subsection{Prime Index Alternatives}

\MillerRabin{} may be substituted with:
\begin{itemize}[noitemsep]
    \item \textbf{AKS:} Deterministic primality test
    \item \textbf{Precomputed prime tables:} For bounded domains
\end{itemize}

\subsection{Alternative Drift Metrics}

The drift bound $\delta(t) \leq 0.3$ may be replaced with domain-specific metrics (e.g., Wasserstein distance, KL divergence) provided the contractive property (\ref{eq:contractive}) is preserved.

\subsection{Alternative Execution Environments}

The reference implementation targets EVM-compatible chains but may be adapted to:
\begin{itemize}[noitemsep]
    \item \textbf{WASM runtimes:} (Substrate, Cosmos)
    \item \textbf{TEE environments:} (SGX, TrustZone)
    \item \textbf{Pure off-chain execution:} with periodic anchoring
\end{itemize}

%% ============================================================
\section*{Document Metadata}

\begin{table}[h]
\centering
\begin{tabular}{@{}ll@{}}
\toprule
\textbf{Field} & \textbf{Value} \\ \midrule
Canonical Title & MTPI Core Specification v1.0 \\
Subtitle & A Proof-First Framework for Verifiable Computation \\
Version & 1.0.0 \\
Publication Date & January 14, 2026 \\
Author & Van Gelder, Ryan O. \\
Affiliation & Citizen Gardens Institute for Mathematical Discovery \\
License & CC-BY-4.0 International \\
Document Type & Defensive Publication / Technical Note \\
DOI & \href{https://doi.org/10.5281/zenodo.18248116}{\texttt{10.5281/zenodo.18248116}} \\
SHA256 & \texttt{bc3599255706f375f116124734cae530258518d391dc745e51e9a73ed140e936} \\
\bottomrule
\end{tabular}
\end{table}

%% ============================================================
\begin{thebibliography}{9}

\bibitem{groth2016}
J. Groth, ``On the Size of Pairing-Based Non-interactive Arguments,'' 
\textit{EUROCRYPT 2016}, LNCS 9666, pp. 305--326, 2016.
DOI: 10.1007/978-3-662-49896-5\_5

\bibitem{grassi2021}
L. Grassi, D. Khovratovich, C. Rechberger, A. Roy, and M. Schofnegger,
``Poseidon: A New Hash Function for Zero-Knowledge Proof Systems,''
\textit{USENIX Security Symposium}, 2021.

\bibitem{rabin1980}
M. O. Rabin, ``Probabilistic algorithm for testing primality,''
\textit{Journal of Number Theory}, vol. 12, no. 1, pp. 128--138, 1980.
DOI: 10.1016/0022-314X(80)90084-0

\bibitem{vanosdol2025}
T. Van Osdol, ``The Dynamic Recursive Operator $\Xi(t)$ in PIRTM Systems,''
Internal Framework Notes, TVO Doc 6, 2025.
Defines recursive tensor evolution models using $\Lambda_m$, Hilbert-Schmidt dynamics, 
and prime-indexed contraction mappings. Shows stability in DRMM/PIRTM via entangled feedback loops.

\end{thebibliography}

%% ============================================================
\appendix
\section{Mathematical Notation Reference}

\begin{table}[h]
\centering
\begin{tabular}{@{}ll@{}}
\toprule
\textbf{Symbol} & \textbf{Meaning} \\ \midrule
$\lambda$ & Universal Multiplicity Constant (cognitive stabilizer), $\lambda \geq 1$ \\
$\Tensor{t}$ & Tensor state at time $t$ \\
$\mathbb{P}(N)$ & Set of primes up to index $N$ \\
$\delta(t)$ & Drift metric at time $t$ \\
$[\mathcal{M}, A]$ & Commutator of operators $\mathcal{M}$ and $A$ \\
$\mathcal{E}_t$ & Ethical Tensor Field at time $t$ \\
$\sigma_i(t)$ & Sovereignty Tensor for agent $i$ at time $t$ \\
$H_n$ & Hash chain state at block $n$ \\
$Q_{\text{ent}}(n)$ & High-entropy source at index $n$ \\
$\Phi_t$ & Recursive operator at time $t$ \\
$\Xi(t)$ & Dynamic recursive operator (Van Osdol) \\
$\Lambda_m$ & Prime-indexed contraction mapping parameter \\
\bottomrule
\end{tabular}
\caption{Mathematical notation reference}
\end{table}

%% ============================================================
\section{Conformance Checklist}

An implementation claiming MTPI v1.0.0 conformance \textbf{MUST} satisfy:

\begin{table}[h]
\centering
\begin{tabular}{@{}lll@{}}
\toprule
\textbf{Requirement} & \textbf{Test} & \textbf{Pass Condition} \\ \midrule
Prime-Gated Activation & Miller-Rabin verification & $\text{isPrime}(\text{epochIndex}) = \text{true}$ \\
Drift Bound & Client + on-chain check & $\delta(t) \leq 0.3$ \\
CSL Commutation & Proof verification & $[\mathcal{M}, \mathcal{E}_t] = 0$ for all transitions \\
Replay Protection & Nullifier lookup & $\text{usedProofHashes}[\cdot] = \text{false}$ \\
Silence-by-Default & Failure mode test & Revert/halt on invalid proof \\
Zero-Surveillance & Emission Gating Test & No unlisted outbound connections \\
Client-Side Proving & Architecture audit & Proof generation on user device \\
Archivum Compliance & Data audit & No raw PII in audit logs \\
\bottomrule
\end{tabular}
\caption{Conformance checklist for MTPI v1.0.0}
\end{table}

\vfill

\begin{center}
\rule{0.5\textwidth}{0.4pt}\\[0.5cm]
\textit{A specification that cannot be copied is a secret.\\
A specification that can be copied but not enclosed is prior art.}\\[0.5cm]
\textbf{END OF SPECIFICATION v1.0.0}
\end{center}

\end{document}
